\section{Height one quartics by signed occurrence fusion}
\label{sec:signed-network}
The new reduction acts on a finite quadratic system over unrestricted
integers. It changes neither its integer zero set up to a specified
bijection nor its inherited consistency calibration. Its new ingredients
are row-sign balancing and \emph{opposite-occurrence fusion}: equal
monomials carrying opposite source signs are represented by a single
incidence edge between their source rows. The retained Round~17 incidence compiler routed
both occurrences separately to a common hub; the incoming Graph digits
package also used occurrence pairing during alias identification. Here the
pair itself becomes a direct source-row edge. All claims about coefficient
height below refer to complete collected coefficient files, not merely an
arithmetic circuit.

\subsection{Certified normalization and row orientations}
Let the given source be $R_r(x)=0$. We first use the preceding exact affine
refactorization, when applicable: if the homogeneous quadratic part of a
row is $\eta L(x)^2$, with $\eta=\pm1$, introduce $z=L(x)+k$ and replace
the row by $\eta z^2+R_r(x)-\eta(L(x)+k)^2=0$. The new definition is monic
and its value is uniquely determined. Shared affine definitions can be
reused. Coefficient-by-coefficient substitution verifies every transformed
row. The provided original source, transformed source and factor
certificate fix this operation completely.

Normalize integer coefficients by monic definitions of integer constants
$t_c=c$ and scale variables $s_{c,i}=cX_i$. A constant is defined from
already defined smaller constants as $t_c=t_a+t_b$ when $a+b=c$, using
$t_1=1$. If no such pair is available, introduce the smaller binary
intermediate and use
$t_c=t_{\lfloor c/2\rfloor}+t_{\lfloor c/2\rfloor}+(c\bmod2)$.
Each added definition is triangular and monic. The delivered DAG prefers
two already-needed constants; in particular it does not introduce a
binary intermediate merely because a binary-only recipe would do so.

For each absolute quadratic coefficient $c>1$, form a finite graph on the
source coordinates, with an edge $\{i,j\}$ for every term $cX_iX_j$.
A loop at $i$ requires selecting $i$. Choose any vertex cover $V_c$ of
this graph, and introduce $s_{c,i}=t_cX_i$ only for $i\in V_c$. A term
$cX_iX_j$ is then represented by $s_{c,i}X_j$ if $i$ was selected, or by
$X_is_{c,j}$ otherwise. The cover condition proves that one option always
exists. The supplied deterministic greedy cover starts with the loop
vertices and repeatedly selects a highest-degree remaining vertex;
optimality is not assumed. Constants and linear terms use $t_c$ and
$t_cX_i$, respectively. Unit coefficients need no constructor.

Splitting repeated summands into separate occurrences yields
\[
 S_r(y)=\sum_j\epsilon_{rj}m_{rj}(y),\qquad \epsilon_{rj}\in\{-1,1\},
 \quad \deg m_{rj}\le2.
\]
Every normalized coordinate has a unique integral value over an original
tuple. The checker verifies all triangular pivots and each exact source-row
pullback, rather than trusting the constant-chain or cover-selection
recipe. Thus this refinement saves constructors without changing either
witness map. The coefficient-factor choice, just like the later graph,
is fixed by the certificate and is not a parameter of the polynomial.

Choose a fixed sign $\sigma_r\in\{-1,1\}$ for each normalized row, and
replace $S_r$ by $T_r=\sigma_rS_r$. This preserves both its zero set and
uniqueness of all constructors. The choice is part of the finite
certificate, not a parameter of the final polynomial. The supplied
heuristic flips a row when doing so reduces
\[
 \sum_m\left|\sum_{r,j:m_{rj}=m}\sigma_r\epsilon_{rj}\right|.
\]
This balances opposite occurrences of the same logical monomial.
Correctness does not require optimality of this heuristic.

\subsection{Direct edges for cancelling occurrences}
Create disjoint vertex blocks $B_r$ for the rows and a hub block $B_0$.
Inside each block choose a tree, with one independent unrestricted integer
flow variable on each tree edge. A normalized occurrence is assigned to a
vertex in its own row block. Partition the occurrences into sets of size
one or two. A pair is allowed only when its occurrences are in distinct
rows, have the same logical monomial, and have opposite signs in $T_r$.

A singleton of sign $\epsilon$ is represented by an edge from its row
vertex to a hub vertex carrying $\epsilon M$. A pair $+m$ in row $r$ and
$-m$ in row $s$ is represented by one edge from its row-$r$ vertex to its
row-$s$ vertex, carrying $M$. Here $M$ is a degree-two homogeneous
monomial in aliases of the logical factors of $m$: a linear term is
multiplied by a unit alias and a constant uses two unit aliases. Each factor
alias has one declared logical value. The two occurrences of a pair use
\emph{the same} edge and the same factor aliases.

The graph is required to be simple. The constructor rejects a pairing if
its proposed unordered vertex pair was already used; remaining hub edges
are placed without parallel edges. The certificate lists every bound
source occurrence, and the checker verifies that each appears exactly
once. More precisely, the edge boundary at the row-block level equals its
list of signed source occurrences. Consequently, after substituting the
intended logical values,
\begin{equation}\label{eq:networkboundary}
 \sum_{v\in B_r} b_v=T_r(y),\qquad
 \sum_{v\in B_0} b_v=-\sum_rT_r(y).
\end{equation}
where $b_v$ is the divergence from prescribed edges alone. A fused edge
contributes $+M$ and $-M$ to precisely its two bound rows; hence \eqref{eq:networkboundary} is
unchanged by fusion. This is the central new identity.

\subsection{The coefficient-height-one energy}
Let $d_v$ denote the outgoing-minus-incoming sum of all prescribed and
free edges, and put
\[
 E=\frac12\sum_vd_v^2.
\]
Each edge has two endpoints, so its diagonal energy coefficient is one.
Two distinct edges of a simple graph meet at most once, so their formal
cross coefficient is $0$ or $\pm1$. Give each prescribed product an
initial private factor alias, and share its other logical factor as in the
previous incidence compiler. The initial fully expanded energy has height
one. Subsequently identify only aliases of the same logical value, and
accept an identification only when full exact coefficient collection
still has coefficients in $\{-1,0,1\}$. Such substitution preserves the
sum-of-squares identity even when signed collisions cancel.

For every remaining non-unit logical class with at least three aliases use
\[
 C(a_1,\ldots,a_k)=\sum_i a_i^2-\sum_i a_i a_{i+1},\quad a_{k+1}=a_1.
\]
Since $2C=\sum_i(a_i-a_{i+1})^2$, it is nonnegative and vanishes exactly
when all aliases agree. A singleton needs no penalty. A class with two
aliases is either completed to a three-cycle by one dummy, or uses
\[
 Q(a,b)=a^4+b^4-a^3b-ab^3=(a-b)^2(a^2+ab+b^2).
\]
This is nonnegative and vanishes exactly at $a=b$; also
$2Q=(a^2-b^2)^2+(a^2-ab)^2+(ab-b^2)^2$. The latter option is used only
when neither alias occurs as a prescribed square, which ensures its four
monomials cannot collide with the incidence energy. For all unit aliases,
adding dummies if necessary so that $\ell\ge2$, put
\[
 U=1+\sum_{i=1}^{\ell}u_i^2-u_1-u_\ell-
       \sum_{i=1}^{\ell-1}u_iu_{i+1}.
\]
Twice $U$ is $(u_1-1)^2+\sum_{i<\ell}(u_i-u_{i+1})^2+(u_\ell-1)^2$.
Thus the final literal polynomial is
\begin{equation}\label{eq:networkenergy}
 \boxed{P=E+\sum C+\sum Q+U.}
\end{equation}
It is nonnegative over the reals. All monomials have total degree at most
four. The full coefficient files are independently checked to have exact
degree four, every nonzero coefficient equal to $+1$ or $-1$, constant
coefficient one, and every stated variable present. The new fusion step
changes the source-edge incidence graph, but requires no change to these
equality penalties or their proof.

\subsection{Forward map, inverse map, and uniqueness}
Given a source tuple, evaluate the monic affine, constant and scale
constructors, assign each ordinary alias its logical value and every unit
alias the value one. Assign equality dummies the value of their class.
All equality and unit penalties then vanish.

Root each internal tree at its first listed vertex. For a tree edge with
child $v$, let $\tau_v=\pm1$ be its incidence sign at the child, and set
\[
 w_v=-\tau_v\sum_{t\text{ in the subtree of }v}b_t.
\]
This is an explicit integral polynomial in the source coordinates.
Eliminating leaves proves that all nonroot divergences vanish. Formula
\eqref{eq:networkboundary} gives root divergence $T_r$ in each row block and $-\sum_rT_r$ in the
hub block. Thus the canonical lift satisfies, universally,
\begin{equation}\label{eq:networklift}
 P(L(x))=\frac12\left(\sum_rR_r(x)^2+
                \left(\sum_r\delta_rR_r(x)\right)^2\right).
\end{equation}
where $\delta_r$ is the certified sign of the normalized row descending
from the original residual $R_r$. Constructor rows pull back to zero.
In particular every source zero has an integer lift. The affine-definition
rows also vanish identically under their canonical substitution, so \eqref{eq:networklift}
composes unchanged through affine refactorization. One must not replace
\eqref{eq:networklift} by a bare sum of the original squares.

Conversely, if $P(z)=0$, nonnegativity in \eqref{eq:networkenergy} forces every equality
penalty and $U$ to vanish. Hence all ordinary aliases in a class agree,
and all unit aliases equal one. Also $E=0$ forces every vertex divergence
to vanish. Summing them over a row block cancels its internal tree flows;
by the certified occurrence-boundary identity the result is precisely
$T_r(y)=0$. Since $\sigma_r=\pm1$, all normalized rows $S_r(y)$ vanish.
Triangular monic constructors then force the constant, scale and affine
coordinates to their intended values. Reading the original logical
coordinates yields an original source zero.

The same argument proves uniqueness: constructors, aliases, units and
dummies have unique values over the original tuple, and leaf elimination
uniquely determines every internal tree flow. Thus the displayed forward
map and coordinate projection are inverse bijections between integer zero
sets. No sign restriction on a final variable, division by a variable,
or unproved boundedness assertion is used.

\subsection{Evidence and inherited consistency scope}
The independent checker imports no producer. It verifies source hashes,
all signed row pullbacks, topological constructors, every occurrence and
its edge boundary, graph simplicity, every internal tree, equality-class
coverage, all coefficients through literal half-sum-of-squares expansion,
and every vertex of the universal symbolic lift. Its tests use explicit
exceptions and remain active under optimized Python. The complete
polynomial is also supplied in ordinary literal notation. Neither this
finite algebra nor its successful replay proves $\operatorname{Con}(\mathrm{ZFC})$:
it transfers the supplied source's existing ZFC-provable equivalence
between unsolvability and consistency. No global-optimality, literature-wide
priority or new end-to-end formal-kernel calibration is claimed.
